\section{撤藩：把临时安排变成战争}

三藩并非清初版图上一块静止的颜色。平西、平南、靖南三藩拥有驻军、地盘和饷源，是征服西南和两广的工具，也使中央财政与皇权面临一个不能无限期延宕的问题。康熙十二年三月，朝廷批准或安排撤藩（EQ-1673-FEUDATORY-ABOLITION）。《清圣祖实录》、内阁大库题本与《平定三逆方略》能确立决策的文书序列；各藩请求的先后、朝议的立场以及军费核算，仍须把奏报、部议与最终处分拆开阅读。撤藩不是一句抽象的“削藩”，而是把将领家族的继承、军队的粮饷和地方中介的前途同时置于风险之中。

吴三桂于同年十一月在云南起兵（EQ-1673-WU-REVOLT），耿精忠次年在福建反叛，王辅臣据平凉起事（EQ-1674-GENG-JINGZHONG、EQ-1674-WANG-FUCHEN）。战争迅速把云南、福建和西北交通连在一起。吴军檄文中的忠明、族群和财政语言是动员资源，不可直接等同所有参与者的动机；《平定三逆方略》是清廷战时汇编，也会按凯旋叙事组织敌我。福建的耿精忠既需郑氏海上力量，又要维持陆上府县，郑经同年自台湾入闽、占据若干沿海节点（EQ-1674-ZHENG-JING-FUJIAN）。据点易手可由军报、方志和郑氏材料相互钉住，盟约的范围、征粮和实际控制却只能逐地核查。

两广的尚之信在1675年前后控制广州军政、反复周旋（EQ-1675-SHANG-ZHIXIN）。这类“反清”或“归清”的标签若不按命令、行动和时间分段，就会遮蔽港口财政、家族继承与地方军队的讨价还价。1676年耿精忠投降，清军收复武昌等长江中游节点（EQ-1676-GENG-SURRENDER、EQ-1676-WUCHANG）。赦降文书说明招抚的条件，不能自动说明降军已被安置、郑氏影响已消退；一座府城的收复，也不等于周边县乡已脱离战争。

吴三桂1677年在衡阳称帝、国号周，试图以更高的政治名义整合将领、赋役和盟友（EQ-1677-WU-ZHOU）。称帝本身可靠，礼制细节和各地响应范围却主要见于敌对军报，需同吴周檄文互证。翌年吴三桂去世，吴世璠继承（EQ-1678-WU-DEATH）；此后清军在湖南推进、清理四川通往西南的道路（EQ-1679-HUNAN-RECOVERY、EQ-1680-SICHUAN-RECOVERY）。《实录》与方略记录了军事推进，地方志、军需题本和契约材料才可能局部说明逃亡、仓储、垦复和赋役恢复。特别是四川战后人口与耕地数字，口径差异很大，不应从任何一种招民叙事中推出整齐的复苏。

康熙二十年十月，清军入昆明，吴世璠自尽，吴周政权结束（EQ-1681-KUNMING）。城陷与死亡是可靠的终局节点；围城损失、降附与后续清剿并非同日完成。多年战争中形成的驻军、饷运、道路损耗与地方恢复，仍将塑造云南、贵州和两广的治理。中央收回的是名义军政权，却继承了把这一权力落到山地、城镇和税源上的长期成本。
